\subsection{Analyse der HTTP/2-Kommunikation}

Die HTTP/2-Kommunikation zwischend dem Client und dem Server wurde über Wireshark analysiert. Dabei wurden die folgenden Schritte durchgeführt:

\begin{enumerate}
    \item Starten von Wireshark und Auswahl der entsprechenden Netzwerkschnittstelle, um den Datenverkehr zu erfassen.
    \item Anfrage an die Webseite mit dem Python Skript \ref{code:connection} senden, um die HTTP/2-Kommunikation zu initiieren.
    \subitem \texttt{python3 connection.py --http\_version 2}
    \item Aufgezeichnete Pakete nach der Anfrage durchsuchen.
    \item Filterung der Frame Pakete Nummern um nur DNS Abfrage, TCP Handshake und HTTP/2 Pakete zu betrachten.
    \subitem \texttt{frame.number >= \{ErstesPaketNummer\} \&\& frame.number <= \{LetztesPaketNummer\}}
    \item Speichern der relevanten Pakete für die weitere Analyse.
    \item Analyse der HTTP/2 Pakete, um die Kommunikation zwischen Client und Server zu verstehen, einschließlich der Anfragen und Antworten, die über HTTP/2 gesendet wurden.
\end{enumerate}

